\documentclass[12pt]{article}
\usepackage{amsmath}
\usepackage{amssymb}
\usepackage{physics}
\usepackage{graphicx}
\usepackage{hyperref}
\usepackage{booktabs}
\usepackage{array}

\newcommand{\ip}[2]{\langle #1 \mid #2 \rangle}

\title{The Transformation of Time and Space}
\author{}
\date{}

\begin{document}
\maketitle
 
% ============================================================
\section{Introduction}
% ============================================================

In his seminal article, Einstein elegantly derived the Lorentz transformation for time and
space between systems moving at constant relative velocity. His original derivation assumed
that the relative velocity was aligned with the $x$-axis, and implicitly considered the
measurement of length perpendicular to this velocity direction to be consistent in both
systems. In this section, we will follow Einstein's reasoning and extend the derivation to a
more general case---a boost in any arbitrary direction. Our goals are as follows:

\begin{enumerate}
  \item \textbf{General Boost Transformation:}
    \begin{itemize}
      \item We will derive the Lorentz transformation for a general boost, where the relative
            velocity is represented by an arbitrary vector.
      \item This broader approach allows us to handle boosts in any direction, not just along
            the $x$-axis.
    \end{itemize}
  \item \textbf{Preserving Length Perpendicular to Velocity:}
    \begin{itemize}
      \item We will explicitly discuss the conditions that ensure the preservation of lengths
            perpendicular to the velocity direction.
      \item This consideration is crucial for maintaining consistency across reference frames.
    \end{itemize}
  \item \textbf{Time Transformation via Clock Synchronization:}
    \begin{itemize}
      \item Following Einstein's original method, we will derive the transformation of time
            based on clock synchronization.
      \item This step lays the foundation for our subsequent analysis.
    \end{itemize}
  \item \textbf{Lorentz Transformation for Time and Space:}
    \begin{itemize}
      \item Leveraging the invariance of the speed of light, we will derive a concise
            expression for both time and space transformations.
    \end{itemize}
\end{enumerate}

\textbf{Assumptions and Prerequisites:} We assume that readers are familiar with advanced
calculus, linear algebra, and the fundamental principles of Special Relativity. We will make
use of the bra-ket notation, which readers can get familiar with at
\url{https://en.wikipedia.org/wiki/Bra\%E2\%80\%93ket_notation}.

% ============================================================
\section{Synchronization of Clocks}
% ============================================================

In Special Relativity, the concept of clock synchronization assumes a pivotal role in
developing the Lorentz transformation of time and space. Although seemingly straightforward,
we will delve into a detailed explanation of this fundamental concept, as defining
synchronization becomes essential for determining when events occur simultaneously.

Let us delve into the starting point of our exploration. Imagine two observers, A and B, both
at rest within the same reference system but positioned at different locations. These
observers, A and B, each have their own clocks, measuring time independently.

Now, let us imagine an experiment where a photon is emitted from A, travels to B, and gets
reflected back to its origin in A. Observer A records the time of emission as $t_1$, while
observer B notes the time of reflection as $t_2$. Finally, A marks the time of the photon
arrival as $t_3$. We assume that both A and B have access to each other's time recordings
(for example they would exchange data after the experiment). A possible outcome of this
experiment could be as per the table below (units of time are not relevant in this
discussion):

\begin{center}
\begin{tabular}{ccc}
\toprule
Time in A & Time in B & Event \\
\midrule
1 & 2 & Emission, $t_1 = 1$ \\
3 & 4 & Reflection, $t_2 = 4$ \\
5 & 6 & Arrival, $t_3 = 5$ \\
\bottomrule
\end{tabular}
\end{center}

If the observer at point B used the time difference $t_2 - t_1$ to compute the time needed
for the photon to reach B, they would obtain an incorrect result of 3 instead of the actual
value of 2. More worrisome, the time difference $(t_3 - t_2)$ would be 1, indicating a
relative motion of the two observers who are at rest in the same reference system. These
discrepancies occur because the two clocks are not synchronized.

The way for B to synchronize their clock with A is to record the time stamp $t_1$, wait for
the photon to get back to A, ask A for $t_3$, and set the clock in B such that:
\begin{equation}
  t_2 - t_1 = t_3 - t_2
  \label{eq:clock_synchronization}
\end{equation}
which is realized in this case by setting their clock back 1 unit.

The synchronization of clocks is essential for determining how time and space are measured
when we consider events observed by two systems of reference now in motion with respect to
each other. This is the subject of the next sections.

% ============================================================
\section{Transformation of Time}
% ============================================================

In this section we will derive the equation for the transformation of time between two
systems of reference moving at constant relative velocity.

\subsection{The Mirror Experiment}

Let us explore the mirror experiment from the perspective of a new reference frame. In this
scenario, the mirror and the emitting source are both at rest within a system of reference
$F'$, which itself is in motion relative to another system of reference $F$. Observers in $F$
measure time and distance as $t$ and $x_1, x_2, x_3$, while observers in $F'$ measure time
and distance as $t'$ and $x'_1, x'_2, x'_3$. We assume that at time zero the origin of the
systems coincide.

\subsubsection*{Experiment Observed from $F'$}

The source is located at the origin of $F'$ and the mirror is positioned at a fixed distance
along an arbitrary direction. A photon is emitted at a time $t'_0$, it reaches the mirror at
time $t'_1$ where it is reflected, and it arrives back at the source at time $t'_2$.

\begin{figure}[h]
  \centering
  \includegraphics[width=0.6\textwidth]{movingframes.jpg}
  \caption{Fig~1: Photon Emission and Reflection in $F'$}
\end{figure}

As we assume that the clocks in $F'$ are synchronized, the following relation holds true:
\begin{equation}
  \frac{1}{2}(t'_0 + t'_2) = t'_1
  \label{eq:clock_synchronization_Fprime}
\end{equation}

\subsubsection*{Experiment Observed from $F$}

Now we examine the three events---emission, reflection, and arrival back at the source---from
the viewpoint of reference frame $F$. It is crucial to emphasize that all time, velocity, and
distance quantities in this section \textit{are assumed to be measured within the system of
reference} $F$.

The photon is emitted from the source at time $t_0$ and at the position $\vb{v}\, t_0$ where
$\vb{v}$ is the velocity of $F'$ measured from $F$.

\begin{figure}[h]
  \centering
  \includegraphics[width=0.6\textwidth]{movingframes1.jpg}
  \caption{Fig~2: Velocities and positions of the photon seen from $F$}
\end{figure}

The photon travels to the mirror with speed $\vb{c}_{up} = \vb{c}_1 + \vb{v}$, and it
reaches it at a time:
\begin{equation}
  t_1 = t_0 + \frac{l}{|\vb{c}_1|} = t_0 + \frac{l}{|\vb{c}_{up} - \vb{v}|}
\end{equation}
where $l$ is the distance of the mirror from the source measured from $F$, and $\vb{c}_1$ is
the photon velocity in $F'$ as observed in $F$. When the photon reaches the mirror its
position will be
\begin{equation}
  \vb{v}\, t_1 + (\vb{c}_{up} - \vb{v})(t_1 - t_0)
  = \vb{v}\, t_0 + l\,\frac{\vb{c}_{up}}{|\vb{c}_{up} - \vb{v}|}
\end{equation}

The time $t_2$ when the photon reaches the source back again is
\begin{equation}
  t_2 = t_1 + \frac{l}{|\vb{c}_{down} - \vb{v}|}
      = t_0 + \frac{l}{|\vb{c}_{up} - \vb{v}|} + \frac{l}{|\vb{c}_{down} - \vb{v}|}
\end{equation}
where $\vb{c}_{down}$ is the photon velocity returning to the source. The position of the
photon when it meets the source is,
\begin{equation}
  \vb{v}\, t_2 = \vb{v}\, t_0
    + \vb{v}\, \frac{l}{|\vb{c}_{up} - \vb{v}|}
    + \frac{l}{|\vb{c}_{down} - \vb{v}|}.
\end{equation}

To derive the Lorentz transformation, we will assume the existence of a functional
relationship between $t'$ and $t$,
\begin{align}
  t' &: F \rightarrow \mathbb{R} \\
  t' &= t'(\vb{x}, t)
\end{align}
and formally, the events of emission, reflection and meeting the source again will be observed
in $F'$ as
\begin{align}
  t'_0 &= t'(\vb{v}\, t_0,\; t_0) \\
  t'_1 &= t'\!\left(
    \vb{v}\, t_0 + l\,\frac{\vb{c}_{up}}{|\vb{c}_{up} - \vb{v}|},\;
    t_0 + \frac{l}{|\vb{c}_{up} - \vb{v}|}
  \right) \\
  t'_2 &= t'\!\left(
    \vb{v}\, t_0 + \vb{v}\, l
      \bigl(|\vb{c}_{up} - \vb{v}|^{-1} + |\vb{c}_{down} - \vb{v}|^{-1}\bigr),\;
    t_0 + l\bigl(|\vb{c}_{up} - \vb{v}|^{-1} + |\vb{c}_{down} - \vb{v}|^{-1}\bigr)
  \right)
\end{align}

Expanding the function $f(l): \mathbb{R} \rightarrow \mathbb{R}$,
\begin{equation}
  f(l) = t'\!\left(
    \vb{v}\, t_0 + l\,\frac{\vb{c}_{up}}{|\vb{c}_{up} - \vb{v}|},\;
    t_0 + \frac{l}{|\vb{c}_{up} - \vb{v}|}
  \right)
\end{equation}
in Taylor series around zero we obtain,
\begin{equation}
  t'_1 = t'(\vb{v}\, t_0,\; t_0)
    + l\,|\vb{c}_{up} - \vb{v}|^{-1}
      \ip{\grad t'}{(\vb{c}_{up}, 1)} + o(l^2).
\end{equation}

A similar expression holds true for $t'_2$,
\begin{equation}
  t'_2 = t'(\vb{v}\, t_0)
    + l\bigl(|\vb{c}_{up} - \vb{v}|^{-1} + |\vb{c}_{down} - \vb{v}|^{-1}\bigr)
      \ip{\grad t'}{(\vb{v}, 1)} + o(l^2)
\end{equation}

Using the synchronization condition of \eqref{eq:clock_synchronization_Fprime}, we arrive at
the differential equation in $t'$ (for $l \to 0$),
\begin{equation}
  \ip{\grad t'}{\vb{w}} = 0
\end{equation}
where,
\begin{equation}
  \vb{w} = \Bigl(
    \vb{v}\bigl(|\vb{c}_{up} - \vb{v}|^{-1} + |\vb{c}_{down} - \vb{v}|^{-1}\bigr)
    - 2\vb{c}_{up}|\vb{c}_{up} - \vb{v}|^{-1},\;
    |\vb{c}_{up} - \vb{v}|^{-1} - |\vb{c}_{down} - \vb{v}|^{-1}
  \Bigr).
\end{equation}

Assuming a linear relationship between the time in $F'$ and the coordinates in $F$ as,
\begin{equation}
  t' = \vb{p} \cdot \vb{x} + \beta\, t
\end{equation}
where $\vb{p}$ is a vector transforming the spatial coordinates and $\beta$ is the coefficient
transforming time, we arrive at the linear equation in $\vb{p}$ and $\beta$
\begin{equation}
  \vb{p} \cdot \bigl[
    \vb{v}(|\vb{c}_{up} - \vb{v}|^{-1} + |\vb{c}_{down} - \vb{v}|^{-1})
    - 2\vb{c}_{up}|\vb{c}_{up} - \vb{v}|^{-1}
  \bigr]
  + \beta\bigl(|\vb{c}_{up} - \vb{v}|^{-1} - |\vb{c}_{down} - \vb{v}|^{-1}\bigr) = 0
  \label{eq:condition_on_p}
\end{equation}

\begin{figure}[h]
  \centering
  \includegraphics[width=0.5\textwidth]{fig1.png}
  \caption{Fig~3: The photon velocities}
\end{figure}

Equation~\eqref{eq:condition_on_p} can be further simplified if we look in detail at the
direction of the vectors $\vb{v}$, $\vb{c}_{down}$ and $\vb{c}_{up}$. Since the photon in
$F'$ follows the same direction going to the mirror and back to the source,
$\vb{c}_{down} - \vb{v}$ and $\vb{c}_{up} - \vb{v}$ are parallel, but have different
magnitude; hence $\vb{c}_{up}$, $\vb{c}_{down}$ and $\vb{v}$ are all in the same plane, see
Fig.~3.

In the system of reference of three mutually orthogonal unit vectors $\vb{e}_1$, $\vb{e}_2$,
and $\vb{e}_3$, with $\vb{e}_1$ and $\vb{e}_2$ in the same plane of $\vb{c}_{up}$ and
$\vb{v}$, and $\vb{e}_1$ in the same direction of $\mathbf{v}$, the vector $\mathbf{p}$ has
components,
\begin{equation}
  \vb{p} = p_1\,\vb{e}_1 + p_2\,\vb{e}_2 + p_3\,\vb{e}_3
\end{equation}
and equation~\eqref{eq:condition_on_p} simplifies to
\begin{equation}
  p_1\, v\bigl(|\vb{c}_{up} - \vb{v}|^{-1} + |\vb{c}_{down} - \vb{v}|^{-1}\bigr)
  - 2c(p_1\cos\alpha + p_2\sin\alpha)|\vb{c}_{up} - \vb{v}|^{-1}
  + \beta\bigl(|\vb{c}_{up} - \vb{v}|^{-1} - |\vb{c}_{down} - \vb{v}|^{-1}\bigr) = 0
\end{equation}
as $p_3$ can be set to zero as it is orthogonal to the plane of $\vb{c}_{down}$ and $\vb{v}$.
From
\begin{equation}
  \cos\alpha = \frac{c^2 + v^2 - |\vb{c}_{up} - \vb{v}|^2}{2vc}
\end{equation}
we arrive at
\begin{equation}
  p_1\,\frac{c^2}{v} + \beta
  + p_2\,\frac{2c\sin\alpha\,|\vb{c}_{up} - \vb{v}|}{|\vb{c}_{up} - \vb{v}| - |\vb{c}_{down} - \vb{v}|}
  = 0 \quad \forall\, \alpha, v
\end{equation}
which has solutions
\begin{align}
  p_1 &= -\beta\,\frac{v}{c^2} \\
  p_2 &= 0
\end{align}

The transformation of time between $F'$ and $F$ becomes then
\begin{equation}
  ct' = \beta\!\left(-\frac{v}{c}\ip{\vb{e}_1}{\vb{x}} + ct\right)
  \label{eq:transformation_of_time}
\end{equation}
where $\beta$ is a constant depending on the velocity.

% ============================================================
\section{Transformation of Space}
% ============================================================

\subsection{General Form of the Lorentz Transformation}

The Lorentz transformation describes how the coordinates and the time of the same event
observed in $F$ and $F'$ are related to each other, and it is assumed to be linear
\begin{equation}
  \vb{X}' = \boldsymbol{\Lambda}(\vb{v})\,\vb{X}
\end{equation}
where $\boldsymbol{\Lambda}(\vb{v})$ is a matrix depending only on the velocity vector
$\vb{v}$, and $\vb{X}$ is the so-called four-vector
\begin{equation}
  \vb{X} = (x_1, x_2, x_3, ct).
\end{equation}
A four-vector can be represented as a sum of ket vectors formed from a unitary triplet of
orthogonal vectors $\vb{e}_1$, $\vb{e}_2$ and $\vb{e}_3$, with $\vb{e}_1$ chosen parallel to
$\vb{v}$:
\begin{equation}
  \vb{X} = x_1\ket{\vb{e}_1, 0} + x_2\ket{\vb{e}_2, 0} + x_3\ket{\vb{e}_3, 0}
           + ct\ket{\vb{0}, 1}
\end{equation}
and the Lorentz transformation $\boldsymbol{\Lambda}(\vb{v})$ in its most general form
becomes
\begin{equation}
  \begin{aligned}
    \boldsymbol{\Lambda}(\vb{v}) =\;
      &\lambda_{1,1}\ket{\vb{e}_1,0}\bra{\vb{e}_1,0}
      + \lambda_{1,2}\ket{\vb{e}_1,0}\bra{\vb{e}_2,0}
      + \lambda_{1,3}\ket{\vb{e}_1,0}\bra{\vb{e}_3,0}
      + \lambda_{1,4}\ket{\vb{e}_1,0}\bra{\vb{0},1} \\
    +\;&\lambda_{2,1}\ket{\vb{e}_2,0}\bra{\vb{e}_1,0}
      + \lambda_{2,2}\ket{\vb{e}_2,0}\bra{\vb{e}_2,0}
      + \lambda_{2,3}\ket{\vb{e}_2,0}\bra{\vb{e}_3,0}
      + \lambda_{2,4}\ket{\vb{e}_2,0}\bra{\vb{0},1} \\
    +\;&\lambda_{3,1}\ket{\vb{e}_3,0}\bra{\vb{e}_1,0}
      + \lambda_{3,2}\ket{\vb{e}_3,0}\bra{\vb{e}_2,0}
      + \lambda_{3,3}\ket{\vb{e}_3,0}\bra{\vb{e}_3,0}
      + \lambda_{3,4}\ket{\vb{e}_3,0}\bra{\vb{0},1} \\
    +\;&\lambda_{4,1}\ket{\vb{0},1}\bra{\vb{e}_1,0}
      + \lambda_{4,2}\ket{\vb{0},1}\bra{\vb{e}_2,0}
      + \lambda_{4,3}\ket{\vb{0},1}\bra{\vb{e}_3,0}
      + \lambda_{4,4}\ket{\vb{0},1}\bra{\vb{0},1}
  \end{aligned}
  \label{eq:general_Lorentz}
\end{equation}
where $\lambda_{i,j}$ are sixteen free parameters. The transformation of time
\eqref{eq:transformation_of_time} can be utilized to identify some of the $\lambda_{i,j}$ but
not all of them: assumptions need to be made on how events are observed in $F$ and $F'$ to
restrict the possible choices for $\lambda_{i,j}$.

\subsection{The Transformation of Time}

The transformation $\boldsymbol{\Lambda}(\vb{v})$ applied to a four-vector with only a time
component $\vb{X} = x_4\ket{\vb{0},1}$ gives,
\begin{equation}
  \boldsymbol{\Lambda}\, x_4\ket{\vb{0},1}
  = x_4\!\left(
      \lambda_{1,4}\ket{\vb{e}_1,0}
    + \lambda_{2,4}\ket{\vb{e}_2,0}
    + \lambda_{3,4}\ket{\vb{e}_3,0}
    + \lambda_{4,4}\ket{\vb{0},1}
  \right)
\end{equation}
but from the transformation of time \eqref{eq:transformation_of_time} there are only
components along $\vb{e}_1$ and time, hence $\lambda_{2,4} = \lambda_{3,4} = 0$.

\subsection{First Assumption: Preserve Direction Along $\vb{e}_1$}

The first assumption we make is that the direction $\vb{e}_1$ is invariant, meaning that an
event in $F$ with components only along $\vb{e}_1$ will not have components along the other
directions in $F'$. Given an event
$\vb{X} = x_1\ket{\vb{e}_1, 0} + x_4\ket{\vb{0}, 1}$, its transformation in $F'$ will be,
\begin{equation}
  \begin{aligned}
    \boldsymbol{\Lambda}(\vb{v})\bigl(x_1\ket{\vb{e}_1,0} + x_4\ket{\vb{0},1}\bigr)
    &= x_1\!\left(
        \lambda_{1,1}\ket{\vb{e}_1,0}
      + \lambda_{2,1}\ket{\vb{e}_2,0}
      + \lambda_{3,1}\ket{\vb{e}_3,0}
      + \lambda_{4,1}\ket{\vb{0},1}
    \right) \\
    &\quad + x_4\!\left(\lambda_{1,4}\ket{\vb{e}_1,0} + \lambda_{4,4}\ket{\vb{0},1}\right)
  \end{aligned}
\end{equation}
this results in $\lambda_{2,1} = \lambda_{3,1} = 0$ to satisfy our assumption.
The general form \eqref{eq:general_Lorentz} further simplifies as,
\begin{equation}
  \begin{aligned}
    \boldsymbol{\Lambda}(\vb{v}) =\;
      &\lambda_{1,1}\ket{\vb{e}_1,0}\bra{\vb{e}_1,0}
      + \lambda_{1,2}\ket{\vb{e}_1,0}\bra{\vb{e}_2,0}
      + \lambda_{1,3}\ket{\vb{e}_1,0}\bra{\vb{e}_3,0}
      + \lambda_{1,4}\ket{\vb{e}_1,0}\bra{\vb{0},1} \\
    +\;&\lambda_{2,2}\ket{\vb{e}_2,0}\bra{\vb{e}_2,0}
      + \lambda_{2,3}\ket{\vb{e}_2,0}\bra{\vb{e}_3,0} \\
    +\;&\lambda_{3,2}\ket{\vb{e}_3,0}\bra{\vb{e}_2,0}
      + \lambda_{3,3}\ket{\vb{e}_3,0}\bra{\vb{e}_3,0} \\
    +\;&\lambda_{4,1}\ket{\vb{0},1}\bra{\vb{e}_1,0}
      + \lambda_{4,2}\ket{\vb{0},1}\bra{\vb{e}_2,0}
      + \lambda_{4,3}\ket{\vb{0},1}\bra{\vb{e}_3,0}
      + \lambda_{4,4}\ket{\vb{0},1}\bra{\vb{0},1}
  \end{aligned}
\end{equation}

\subsection{Second Assumption: Invariance Orthogonal to $\vb{e}_1$}

The second assumption needed is that the spatial components along $\ket{\vb{e}_2,0}$ and
$\ket{\vb{e}_3,0}$ of any event preserve their direction and length under
$\boldsymbol{\Lambda}(\vb{v})$. For an event
$\vb{X} = x_2\ket{\vb{e}_2, 0} + x_3\ket{\vb{e}_3, 0}$, the corresponding event in $F'$
will be,
\begin{equation}
  \begin{aligned}
    \boldsymbol{\Lambda}(\vb{v})\bigl(x_2\ket{\vb{e}_2,0} + x_3\ket{\vb{e}_3,0}\bigr)
    &= x_2\!\left(
        \lambda_{1,2}\ket{\vb{e}_1,0}
      + \lambda_{2,2}\ket{\vb{e}_2,0}
      + \lambda_{3,2}\ket{\vb{e}_3,0}
    \right) \\
    &\quad + x_3\!\left(
        \lambda_{1,3}\ket{\vb{e}_1,0}
      + \lambda_{2,3}\ket{\vb{e}_2,0}
      + \lambda_{3,3}\ket{\vb{e}_3,0}
    \right).
  \end{aligned}
\end{equation}
Our assumption is satisfied if $\lambda_{1,2} = \lambda_{1,3} = \lambda_{2,3} = \lambda_{3,2}
= 0$ and $\lambda_{2,2} = \lambda_{3,3} = 1$, which results in a further simplification of
$\boldsymbol{\Lambda}(\vb{v})$:
\begin{equation}
  \begin{aligned}
    \boldsymbol{\Lambda}(\vb{v}) =\;
      &\lambda_{1,1}\ket{\vb{e}_1,0}\bra{\vb{e}_1,0}
      + \lambda_{1,4}\ket{\vb{e}_1,0}\bra{\vb{0},1} \\
    +\;&\ket{\vb{e}_2,0}\bra{\vb{e}_2,0} \\
    +\;&\ket{\vb{e}_3,0}\bra{\vb{e}_3,0} \\
    +\;&\lambda_{4,1}\ket{\vb{0},1}\bra{\vb{e}_1,0}
      + \lambda_{4,2}\ket{\vb{0},1}\bra{\vb{e}_2,0}
      + \lambda_{4,3}\ket{\vb{0},1}\bra{\vb{e}_3,0}
      + \lambda_{4,4}\ket{\vb{0},1}\bra{\vb{0},1}
  \end{aligned}
\end{equation}

Using the transformation of time \eqref{eq:transformation_of_time}, we arrive at the final
functional form of $\boldsymbol{\Lambda}(\vb{v})$
\begin{equation}
  \begin{aligned}
    \boldsymbol{\Lambda}(\vb{v}) =\;
      &\lambda_{1,1}\ket{\vb{e}_1,0}\bra{\vb{e}_1,0}
      + \lambda_{1,4}\ket{\vb{e}_1,0}\bra{\vb{0},1}
      + \ket{\vb{e}_2,0}\bra{\vb{e}_2,0}
      + \ket{\vb{e}_3,0}\bra{\vb{e}_3,0} \\
      &- \beta\,\frac{v}{c}\ket{\vb{0},1}\bra{\vb{e}_1,0}
      + \beta\ket{\vb{0},1}\bra{\vb{0},1}
  \end{aligned}
  \label{eq:simplified_Lorentz}
\end{equation}
Finally, the constants $\lambda_{1,1}$, $\lambda_{2,2}$, $\lambda_{3,3}$, and $\beta$ must be
positive to prevent axis reflection.

\subsection{Third Assumption: Invariance of Light Speed}

In Special Relativity the spacetime interval between two events $(x_1, y_1, z_1, ct_1)$ and
$(x_2, y_2, z_2, ct_2)$ is defined as
\begin{equation}
  \Delta s^2 = (x_2 - x_1)^2 + (y_2 - y_1)^2 + (z_2 - z_1)^2 - c^2(t_2 - t_1)^2
\end{equation}
and can be written in bi-linear form as
\begin{equation}
  \Delta s^2 = \ip{\vb{X}}{\vb{M}\,\vb{X}}
\end{equation}
where $\vb{X} = (x_2 - x_1,\; y_2 - y_1,\; z_2 - z_1,\; c(t_2 - t_1))$ and
\begin{equation}
  \vb{M} =
  \begin{pmatrix}
    1 & 0 & 0 &  0 \\
    0 & 1 & 0 &  0 \\
    0 & 0 & 1 &  0 \\
    0 & 0 & 0 & -1
  \end{pmatrix}
\end{equation}

A spacetime interval measured in $F'$ is linked to the same interval in $F$ by the Lorentz
transformation
\begin{equation}
  \Delta s'^2
  = \ip{\vb{X}'}{\vb{M}\,\vb{X}'}
  = \ip{\boldsymbol{\Lambda}\vb{X}}{\vb{M}\,\boldsymbol{\Lambda}\vb{X}}
  = \ip{\vb{X}}{\boldsymbol{\Lambda}^T\vb{M}\,\boldsymbol{\Lambda}\,\vb{X}}
\end{equation}

In the system $F$ we consider a photon moving along a unit vector $\vb{n}$. The photon at
time $t_1$ is in the position $\vb{x}_1$, and moves to the position
$\vb{x}_1 + \vb{n}\,c(t_2 - t_1)$ at the time $t_2$. The spacetime interval in $F$ of the
two events is identically zero,
\begin{equation}
  \Delta s^2 = c^2(t_2 - t_1)^2\|\vb{n}\|^2 - c^2(t_2 - t_1)^2 = 0
\end{equation}

When we examine two events occurring along the trajectory of a photon (referred to as the
photon's world line), the spacetime interval $\Delta' s^2$ in $F'$ must also be zero, assuming
that the speed of light $c$ remains invariant,
\begin{equation}
  \ip{
    \bigl(c(t_2-t_1)\vb{n},\; c(t_2-t_1)\bigr)
  }{
    \boldsymbol{\Lambda}^T\vb{M}\,\boldsymbol{\Lambda}
    \cdot\bigl(c(t_2-t_1)\vb{n},\; c(t_2-t_1)\bigr)
  } = 0
\end{equation}
equivalent to,
\begin{equation}
  \ip{(\vb{n},1)}{\boldsymbol{\Lambda}^T\vb{M}\,\boldsymbol{\Lambda}\cdot(\vb{n},1)}
  = 0
  \label{eq:invariance_light_speed}
\end{equation}
for all $\vb{n}$ such that $\|\vb{n}\| = 1$.

To solve the equation above we need $\vb{M}\,\boldsymbol{\Lambda}$, which is
$\boldsymbol{\Lambda}$ with the time component's sign changed,
\begin{equation}
  \begin{aligned}
    \vb{M}\,\boldsymbol{\Lambda}(\vb{v}) =\;
      &\lambda_{1,1}\ket{\vb{e}_1,0}\bra{\vb{e}_1,0}
      + \lambda_{1,4}\ket{\vb{e}_1,0}\bra{\vb{0},1}
      + \lambda_{2,2}\ket{\vb{e}_2,0}\bra{\vb{e}_2,0}
      + \lambda_{3,3}\ket{\vb{e}_3,0}\bra{\vb{e}_3,0} \\
      &+ \beta\,\frac{v}{c}\ket{\vb{0},1}\bra{\vb{e}_1,0}
      - \beta\ket{\vb{0},1}\bra{\vb{0},1}
  \end{aligned}
\end{equation}

and $\boldsymbol{\Lambda}^T$ is readily obtained by exchanging the bra and ket vectors,
\begin{equation}
  \begin{aligned}
    \boldsymbol{\Lambda}(\vb{v})^T =\;
      &\lambda_{1,1}\ket{\vb{e}_1,0}\bra{\vb{e}_1,0}
      + \lambda_{1,4}\ket{\vb{0},1}\bra{\vb{e}_1,0}
      + \lambda_{2,2}\ket{\vb{e}_2,0}\bra{\vb{e}_2,0}
      + \lambda_{3,3}\ket{\vb{e}_3,0}\bra{\vb{e}_3,0} \\
      &- \beta\,\frac{v}{c}\ket{\vb{e}_1,0}\bra{\vb{0},1}
      + \beta\ket{\vb{0},1}\bra{\vb{0},1}
  \end{aligned}
\end{equation}

The product $\boldsymbol{\Lambda}(\vb{v})^T\vb{M}\,\boldsymbol{\Lambda}(\vb{v})$ is
\begin{equation}
  \begin{aligned}
    &\left(\lambda_{1,1}^2 - \beta^2\frac{v^2}{c^2}\right)\ket{\vb{e}_1,0}\bra{\vb{e}_1,0} \\
    +\;&\left(\lambda_{1,1}\lambda_{1,4} + \beta^2\frac{v}{c}\right)
      \bigl(\ket{\vb{0},1}\bra{\vb{e}_1,0} + \ket{\vb{e}_1,0}\bra{\vb{0},1}\bigr) \\
    +\;&\lambda_{2,2}^2\ket{\vb{e}_2,0}\bra{\vb{e}_2,0}
      + \lambda_{3,3}^2\ket{\vb{e}_3,0}\bra{\vb{e}_3,0} \\
    +\;&\left(\lambda_{1,4}^2 - \beta^2\right)\ket{\vb{0},1}\bra{\vb{0},1}
  \end{aligned}
\end{equation}

and the equation for the invariance of speed of light \eqref{eq:invariance_light_speed} reads,
\begin{equation}
  \left(\lambda_{1,1}^2 - \beta^2\frac{v^2}{c^2}\right)n_1^2
  + 2\left(\lambda_{1,1}\lambda_{1,4} + \beta^2\frac{v}{c}\right)n_1
  + n_2^2 + n_3^2
  + \left(\lambda_{1,4}^2 - \beta^2\right) = 0
\end{equation}

The above equation is identically zero for all $\vb{n}$ such that $\|\vb{n}\| = 1$ only if,
\begin{align}
  \lambda_{1,1}^2 - \beta^2\displaystyle\frac{v^2}{c^2} &= 1 \\
  \lambda_{1,1}\lambda_{1,4} + \beta^2\displaystyle\frac{v}{c} &= 0 \\
  \lambda_{1,4}^2 - \beta^2 &= -1
\end{align}
which has the following admissible solutions,
\begin{align}
  \beta         &= \frac{1}{\displaystyle\sqrt{1 - \frac{v^2}{c^2}}} \\
  \lambda_{1,1} &= \beta \\
  \lambda_{1,4} &= -\frac{v}{c}\,\beta
\end{align}

The Lorentz transformation takes the final form
\begin{equation}
  \boldsymbol{\Lambda}(\vb{v}) =
    \beta\ket{\vb{e}_1,0}\bra{\vb{e}_1,0}
    - \beta\,\frac{v}{c}\ket{\vb{e}_1,0}\bra{\vb{0},1}
    + \ket{\vb{e}_2,0}\bra{\vb{e}_2,0}
    + \ket{\vb{e}_3,0}\bra{\vb{e}_3,0}
    - \beta\,\frac{v}{c}\ket{\vb{0},1}\bra{\vb{e}_1,0}
    + \beta\ket{\vb{0},1}\bra{\vb{0},1}
  \label{eq:Lorentz_transformation}
\end{equation}

The inverse of $\boldsymbol{\Lambda}(\vb{v})$ transforms the coordinates from $F'$ to $F$ and
can be readily obtained by exchanging the velocity sign, as it can be verified directly:
\begin{equation}
  \boldsymbol{\Lambda}^{-1}(\vb{v}) = \boldsymbol{\Lambda}(-\vb{v}) =
    \beta\ket{\vb{e}_1,0}\bra{\vb{e}_1,0}
    + \beta\,\frac{v}{c}\ket{\vb{e}_1,0}\bra{\vb{0},1}
    + \ket{\vb{e}_2,0}\bra{\vb{e}_2,0}
    + \ket{\vb{e}_3,0}\bra{\vb{e}_3,0}
    + \beta\,\frac{v}{c}\ket{\vb{0},1}\bra{\vb{e}_1,0}
    + \beta\ket{\vb{0},1}\bra{\vb{0},1}
  \label{eq:Lorentz_inverse}
\end{equation}

The Lorentz transformation written explicitly in terms of the velocity vector, using
$\ket{\vb{e}_1,0} = \ket{\vb{v},0}/v$, reads,
\begin{equation}
  \boldsymbol{\Lambda}(\vb{v}) =
    \frac{\beta}{v^2}\ket{\vb{v},0}\bra{\vb{v},0}
    - \frac{\beta}{c}\ket{\vb{v},0}\bra{\vb{0},1}
    + \ket{\vb{e}_2,0}\bra{\vb{e}_2,0}
    + \ket{\vb{e}_3,0}\bra{\vb{e}_3,0}
    - \frac{\beta}{c}\ket{\vb{0},1}\bra{\vb{v},0}
    + \beta\ket{\vb{0},1}\bra{\vb{0},1}
  \label{eq:Lorentz_in_velocity}
\end{equation}

When $\vb{v}$ has the direction of the $x$-axis, the transformation simplifies to the original
Einstein formula:
\begin{align}
  x' &= \beta(x - vt)               \nonumber \\
  y' &= y                            \nonumber \\
  z' &= z                            \nonumber \\
  t' &= \beta\!\left(t - \frac{v}{c^2}x\right)
  \label{eq:Lorentz_1d}
\end{align}

Finally, it can readily be proven that the Lorentz transformation satisfies the condition
$\boldsymbol{\Lambda}^T\vb{M}\,\boldsymbol{\Lambda} = \vb{M}$, hence $\Delta' s = \Delta s$
for all events, not only those on a world line of a photon:
\begin{equation}
  \Delta s'^2
  = \ip{\vb{X}'}{\vb{M}\,\vb{X}'}
  = \ip{\boldsymbol{\Lambda}\vb{X}}{\vb{M}\,\boldsymbol{\Lambda}\vb{X}}
  = \ip{\vb{X}}{\boldsymbol{\Lambda}^T\vb{M}\,\boldsymbol{\Lambda}\,\vb{X}}
  = \ip{\vb{X}}{\vb{M}\,\vb{X}}
  = \Delta s
  \label{eq:spacetime_interval_invariance}
\end{equation}

\end{document}
